\subsection{Logarithmic wavefields and possible scaling limits}
\label{S:future-universality}

A logarithmic change of variables provides another way to view random wave
propagation. This idea has classical precedents in the Rytov transformation,
which represents a complex wavefield through its logarithmic amplitude and
phase. For example, Sung et al.~\cite{sung.choi.ea:09:optical}*{Section~2}
substitute $U=e^\psi$ into a Helmholtz equation before applying a
weak-scattering approximation. The substitution and the subsequent
approximation are distinct steps. That frequency-domain setting does not
supply a scaling limit for the time-dependent model considered here.

For the smooth equation $u_{tt}=\Delta u+Vu$, set $h=\log|u|$ on a region where
$u$ is nonzero. The ordinary chain rule gives
\begin{equation}\label{E:future-log-wave}
  h_{tt}+(h_t)^2=\Delta h+|\nabla h|^2+V.
\end{equation}
Thus, the logarithmic wave equation retains both a second time derivative and
a quadratic temporal derivative. In the stochastic model, this calculation
applies first to smooth noise approximations. Passing through zeros or to the
singular-noise limit requires a separate argument; we do not assert a global
logarithmic equation for the canonical solution.

A possible mechanism for an effective first-order time evolution can be seen
formally by writing $h=b(t)+\varphi$, where $b$ is a smooth deterministic
centering. Then
\begin{equation}\label{E:future-centered-log-wave}
  \varphi_{tt}+2b'(t)\varphi_t+(\varphi_t)^2
  =\Delta\varphi+|\nabla\varphi|^2+V-b''(t)-(b'(t))^2.
\end{equation}
If $b'$ approached a positive constant and the other temporal terms on the
left became negligible on a fluctuation scale, the term $2b'\varphi_t$ could
produce an effective parabolic evolution. Neither condition follows from
moment growth. In particular, logarithmic moments need not give the typical
centering, and the long temporal memory of the noise must still be accounted
for. This is a heuristic mechanism to investigate.

The ordinary one-dimensional KPZ comparison concerns the centered logarithm
of the positive stochastic heat solution. Its height, space and time scaling
exponents have ratio $1:2:3$; convergence to the KPZ fixed point is a statement about
random fields and their evolution~\cite{wu:26:kpz}*{Theorems~1.6 and~1.8}. The moment normalization
in the present wave problem controls rare amplification. It does not identify
the typical fluctuation size or spatial correlation length. The cubic
upper-tail prediction in the spatial-white case also gives no evidence by
itself for KPZ behavior: the usual KPZ upper-tail power is $3/2$, while a
cubic power appears in a lower-tail regime~\cite{corwin.ghosal:20:kpz}.
Different rare-event scales do not, by themselves, exclude a common limit in
another centered fluctuation regime.

The open question is whether a specified logarithmic amplitude observable
$H$ admits deterministic functions $b$, $S$ and $L$ for which
\[
  \frac{H(Tt,L(T)x)-b(Tt)}{S(T)}
\]
converges to a nondegenerate random field as $T\to\infty$. A logarithm of a
local spatial $L^2$ amplitude is one possible observable that reduces
sensitivity to isolated pointwise zeros. It can itself vanish, so its domain
must also be specified; it differs from $\log|u|$. The scales and topology must
be specified, and a claim of universality would require stability under
changes of regularization and microscopic covariance details. Temporal
correlations and sensitivity to the initial velocity would help determine
whether a scalar amplitude describes the limiting evolution or whether
additional variables remain necessary. We leave this as an exploratory
question, without prescribing exponents or a new universality class.
